\chapter[Los Espejos Rotos]{Los Espejos Rotos\\ \large La fragmentación del alma y la familia.}

\elegantimage{02_fidelidad_y_familia/web/assets/juxtaposition.png}

\lettrine[lines=3, nindent=0.5em, findent=0.2em]{I}{}brahim vivía para el placer insaciable de los sentidos, viendo a las personas como objetos. Esa energía de desunión quedó impresa en su núcleo espiritual como una Causa de fractura emocional profunda.

\section{El Espejo de las Promesas}
\begin{elegantquote}
\textit{Había una vez un hombre encantador que iba de pueblo en pueblo, enamorando corazones y luego marchándose en la oscuridad, dejando tras de sí solo lágrimas e ilusiones rotas. Creía ser libre como el viento...}

\textit{Sin embargo, el universo tiene una memoria infinita. Años más tarde, o quizá en otra vida, este mismo hombre se encontró sentado frente a una mesa enorme, llena de manjares, pero completamente solo. Ningún amigo le visitaba, ningún amor se quedaba a su lado. Sentía un frío en el alma que ningún fuego podía calentar.}

\textit{Aquella soledad no era un castigo cruel del cielo. Eran las mismas lágrimas que él había dejado caer, ahora convertidas en un mar que lo separaba del mundo. Comprendió entonces que cuando rompía el corazón de otro, en realidad, estaba vaciando el suyo propio. Y al derramar su primera lágrima sincera, el universo, con paciencia infinita, le entregó la primera semilla para por fin aprender a amar.}

\end{elegantquote}

\vspace{1em}
\begin{center}
    \textbf{\Large La lealtad es el cimiento de la paz.}
\end{center}
\vspace{1em}

\section{Análisis Kármico y Traducción}
\begin{wrapfigure}{r}{0.5\textwidth}
  \vspace{-1.5em}
  \begin{center}
  \inlineelegantimage[0.45\textwidth]{02_fidelidad_y_familia/web/assets/pasaje_original.jpg}
  \end{center}
  \vspace{-1.5em}
\end{wrapfigure}
La deslealtad narcisista, que busca la saciedad destruyendo los cimientos afectivos, condena al alma a sufrir el mismo y doloroso desarraigo. Jugar con los sentimientos y la intimidad de los demás nos sitúa, inexorablemente, frente al espejo de la profunda soledad y la desunión del propio linaje, obligándonos a cosechar la frialdad que una vez esparcimos.

